\documentclass[journal]{IEEEtran}

\usepackage[T1]{fontenc}
\usepackage{amsmath,amssymb}
\usepackage{amsthm}
\usepackage{flushend}
\usepackage{microtype}
\usepackage[hyphens]{url}
\usepackage[hidelinks]{hyperref}

% Generated from submission/manuscript-metadata.json.
% Run `make metadata` after changing the canonical JSON.
% Do not edit this file by hand.
\newcommand{\PaperMainTitle}{Policy-to-Purpose-to-Byte Continuity for Post-Quantum Peer-to-Peer Sessions}
\newcommand{\PaperSupplementTitle}{Supplementary Material for ``\PaperMainTitle''}
\newcommand{\PaperSupplementPdfTitle}{Supplementary Material for \PaperMainTitle}
\newcommand{\PaperAuthorName}{Zi'ang Li}
\newcommand{\PaperDegree}{MSc by Research (MScR)}
\newcommand{\PaperAcademicRole}{student}
\newcommand{\PaperAffiliation}{Durham University, Durham, United Kingdom}
\newcommand{\PaperCorrespondingEmail}{2403871950@qq.com}
\newcommand{\PaperAuthorDeclaration}{%
  \PaperAuthorName%
  \thanks{The author is an \PaperDegree{} \PaperAcademicRole{} at \PaperAffiliation{}. E-mail: \PaperCorrespondingEmail.}%
}
\newcommand{\PaperConfigureMainDocument}{%
  \hypersetup{%
    pdftitle={\PaperMainTitle},%
    pdfauthor={\PaperAuthorName},%
    pdfsubject={Independent IEEE TIFS manuscript on policy-bound post-quantum peer-to-peer sessions},%
    pdfkeywords={post-quantum cryptography, hybrid key encapsulation, authenticated key exchange, peer-to-peer security}%
  }%
  \title{\PaperMainTitle}%
  \author{\PaperAuthorDeclaration}%
}
\newcommand{\PaperConfigureSupplementDocument}{%
  \hypersetup{%
    pdftitle={\PaperSupplementPdfTitle},%
    pdfauthor={\PaperAuthorName},%
    pdfsubject={Supplementary material for an independent IEEE TIFS manuscript},%
    pdfkeywords={post-quantum cryptography, peer-to-peer security, supplementary material}%
  }%
  \title{\PaperSupplementTitle}%
  \author{\PaperAuthorDeclaration}%
}

\PaperConfigureMainDocument

\newtheorem{definition}{Definition}
\newcommand{\system}{\textsc{BoundSession}}
\newcommand{\contextv}{\textsf{BoundSessionContextV1}}
\newcommand{\grant}{\textsf{BoundSessionGrant}}

\begin{document}

\maketitle

\begin{abstract}
Post-quantum migration in long-lived peer-to-peer applications crosses signed
policy, hybrid key establishment, authenticated handshake, asynchronous session
ownership, and business-effect authorization. Existing work typically verifies
only a subset of this chain. We propose
\emph{policy-to-purpose-to-byte continuity}: each endpoint accepts an
application effect only if its local authenticated decision authorizes the same
peer-visible policy, purpose, wire cryptographic profile, context, and transcript
as the matching peer run; its exact local owner transaction has committed; and a
purpose-specific effect receipt verifies. \system{} is designed to bind the
purpose and policy projection into context-bound PQ/T key establishment, confirm
the complete transcript, derive purpose-separated channel keys, and require
bilateral authorization readiness before business data reaches a trusted effect
committer. We also define a layered assurance boundary that keeps
computational combiner binding, symbolic protocol analysis, state-machine
checking, implementation conformance, final-binary evidence, and physical-device
measurements distinct. This manuscript is an independent design-stage draft:
the integrated proofs, final-source implementation, and cross-platform
application experiments are deliberately reported as pending rather than
inferred from component artifacts.
\end{abstract}

\begin{IEEEkeywords}
post-quantum cryptography, hybrid key encapsulation, authenticated key exchange,
cryptographic agility, downgrade resistance, formal verification, peer-to-peer
security
\end{IEEEkeywords}

\section{Introduction}
\IEEEPARstart{P}{ost-quantum} migration is often described as replacing a
primitive. In a long-lived peer-to-peer (P2P) product, however, the security
decision passes through policy distribution, provider selection, wire
negotiation, key combination, application authorization, asynchronous ownership,
and evidence generation. A failure at any boundary can invalidate a correct
primitive: an old policy can be replayed, a provider can be re-selected after
verification, a context field can be omitted, a stale connection can publish a
new key, or an application can use a session before its security decision is
committed.

Hybrid PQ/T key establishment combines a post-quantum KEM such as ML-KEM with a
traditional component such as X25519. This hedges uncertainty about either
component, but it does not by itself bind every relevant public value, policy
choice, or application interpretation. Recent binding analyses show that key and
ciphertext serialization choices can matter below the protocol layer
\cite{cdm2024,schmieg2024,chempat2025}. At the other end of the stack, generic
channel-binding and exporter mechanisms deliberately leave application-context
semantics to the caller \cite{rfc5705,rfc9266}. Neither observation answers when
a concrete application is authorized to consume a particular session.

This paper asks one question:

\begin{quote}
How can heterogeneous P2P endpoints ensure that a business capability is bound
to the same authenticated migration policy, application purpose, hybrid
realization, transcript, and current session owner?
\end{quote}

We call this property \emph{policy-to-purpose-to-byte continuity}. Here
``byte'' means the exact canonical bytes consumed by the KEM context, signed
transcript, key schedule, grant, and receipt encoders, together with conformance
evidence for their final implementations. It does not mean business payload,
remote attestation, or a verified source-to-binary compiler. The proposed
protocol, \system{}, treats the application capability---not successful key
derivation---as the acceptance event. A closed purpose type distinguishes, for
example, view-only remote desktop, interactive control, and a particular
content-addressed file transfer. The purpose is bound into the pre-KEM context,
the signed transcript, the key schedule, the exact-owner commit, and the final
application receipt.

\textbf{Intended contributions.} This design-stage manuscript is organized
around three intended contributions, each gated by the evidence ledger shipped
with the draft.

\begin{enumerate}
\item A single policy-to-purpose-to-byte contract that makes local provider refinement,
      downgrade behavior, purpose authorization, session ownership, and
      authorization-before-publication plus effect correspondence one security
      surface.
\item A layered formal argument: a computational binding result is used only as
      a component premise, a new active-attacker protocol model establishes
      agreement and secrecy at the composed message boundary, and a separate
      state model covers monotonic policy and asynchronous publication. We do not
      call conformance evidence a mechanized implementation refinement.
\item One portable protocol core and a source-bound evaluation that measures
      matched cryptographic baselines, adversarial ablations, real non-Apple
      interoperability, and purpose-specific physical-device receipts from the
      final candidate.
\end{enumerate}

\noindent\textbf{Draft status.} No integrated theorem or measurement is claimed
in this version. Numerical results will enter the manuscript only after their
raw inputs and final source identity are selected by one immutable evidence
manifest.

\section{Background and Related Work}

\subsection{Migration and downgrade}
TLS~1.3 provides a mature authenticated channel and binds negotiated parameters
into its transcript \cite{rfc8446}. Fallback signaling can detect selected legacy
downgrades \cite{rfc7507}, while NIST guidance treats crypto agility as an
operational lifecycle rather than an algorithm toggle \cite{nistagility}. Our
focus is narrower than a general secure transport and broader than a fallback
signal: a policy-bound attempt must determine which business capability, if any,
can be released after every subsequent boundary.

\subsection{Hybrid binding and realization}
ML-KEM is standardized in FIPS~203 \cite{fips203}. Binding-security work studies
whether a KEM or combiner commits an accepted key to ciphertext and public-key
inputs \cite{cdm2024,schmieg2024,chempat2025}. X-Wing is an important lean hybrid
design point \cite{xwing}. \system{} does not claim a new primitive merely for
hashing more fields. Its proposed contribution is a composed authorization
invariant whose cryptographic projection uses a binding-safe component.

\subsection{Formal and implementation assurance}
EasyCrypt supports machine-checked computational proofs \cite{easycrypt};
Tamarin and ProVerif support symbolic protocol analysis
\cite{tamarin,proverif}. Verified cryptographic libraries demonstrate stronger
source-level implementation links than cross-language test vectors alone
\cite{hacl,formosa}. We therefore label each assurance layer precisely instead
of treating multiple tools as an end-to-end proof of a binary.

\subsection{Authorization, recovery, and audit}
Complete mediation and least privilege motivate checking authority at every
use; capabilities make object-specific authority explicit
\cite{saltzer1975,miller2003}. Crash-safe verification shows that
authorization-before-publication must survive recovery \cite{fscq2015}, while
secure audit logs offer stronger tamper-resistance than a local record
\cite{schneier1999}. These are foundations, not novelty claims. \system{}
composes them with policy-bound hybrid establishment, exact asynchronous
ownership, and purpose-specific effects in one acceptance invariant.

\section{System and Threat Model}

\subsection{System boundary}
Two previously paired endpoints in one managed policy domain establish a
purpose-specific P2P session. On each endpoint, a separate trusted policy
service (or an equivalently isolated boundary) owns a pinned policy root,
monotonic policy state, the hybrid KEM identity, local owner registry, and
server-side decision entries. Sealed presentation, input, and durable-file
committers are also inside this boundary, together with the transactional store, crypto/identity
ports, and receipt-key service. Effect committers execute or directly observe
only their defined platform effect and cannot accept a caller-supplied success
result. The service authenticates callers with OS IPC credentials
and binds each handle to a caller principal, peer, purpose, local owner, epoch,
expiry, and one-time operation sequence. Application modules, general platform
adapters, discovery, and transports are outside that boundary. They may request
a purpose but cannot construct, transfer, widen, or alter a decision or receipt.

An implementation that runs this boundary in the application's address space is
an explicit honest-host profile and cannot claim resistance to hostile local
modules. Peers agree on a wire cryptographic profile, not on byte-identical local
provider implementations.

The design has one protocol core. Swift, Kotlin, C, and other bindings perform
marshalling and platform I/O; they do not implement independent suite,
fallback, transcript, or ownership state machines.

\subsection{Adversary}
The remote adversary controls the network: it observes, drops, delays, reorders,
duplicates, replays, injects, and concurrently triggers pre-authentication input.
It may compromise selected long-term identity, KEM, or ephemeral keys at times
made explicit by each theorem. An untrusted application module may call public
APIs and race session replacement, but cannot read or execute inside the trusted
policy or effect-committer boundary.

The base claims exclude compromise of the trusted policy/session service,
transactional store, crypto/identity ports, receipt-key service, trusted effect
committer, OS, root, or build host. They also exclude malicious
compilers, universal denial-of-service resistance, remote attestation of peer
policy execution, metadata unlinkability, and rollback of the entire device.
Audit integrity is scoped separately: without an external witness, a local hash
chain is diagnostic evidence, not durable non-repudiation.

\subsection{Security goal}
Let $x$ collect the peer identities, peer session, policy, wire profile, purpose,
context and transcript digests, shared grant and bilateral-ready digests, and a
model-level shared purpose-key-schedule identifier. The last is neither a wire
value nor an endpoint-private key handle. Let $\ell_A=(d_A,e_A,o_A,v_A)$ contain $A$'s local decision,
decision epoch, owner binding, and provider identity, with $\ell_B$ defined
symmetrically. For effect receipt $r$, the
central local-to-peer correspondence is below; $B$ denotes the endpoint that
commits the purpose-defined effect.

\begin{equation}
\begin{split}
\mathsf{Accept}_A(x,\ell_A,r) \Rightarrow {} &
\exists!\,\ell_B:\mathsf{Run}_B(x,\ell_B)\\
&\land\;\mathsf{Permit}_A(x,\ell_A)\\
&\land\;\mathsf{Permit}_B(x,\ell_B)\\
&\land\;\mathsf{Finished}_{A,B}(x)\\
&\land\;\mathsf{OwnerTxn}_A(\ell_A)\\
&\land\;\mathsf{OwnerTxn}_B(\ell_B)\\
&\land\;\mathsf{Enable}_A(x,\ell_A) \prec \mathsf{Publish}_A\\
&\land\;\mathsf{Enable}_B(x,\ell_B)\\
&\land\;\exists!\,\mathsf{EffectCommit}_B(x,r)\\
&\land\;\mathsf{Enable}_B \prec \mathsf{EffectCommit}_B\\
&\land\;\mathsf{EffectCommit}_B \prec \mathsf{ReceiptIssued}_B\\
&\land\;\mathsf{ReceiptOK}_A(x,r).
\end{split}
\label{eq:central-goal}
\end{equation}

The matching peer run agrees on the wire profile, not the private local provider
$v_B$. The complete theorem must state compromise escapes, hybrid-secrecy
premises, effect durability or application semantics, receipt ordering, and the
exact meaning of $\prec$ in the state model. For endpoint $E$, $\mathsf{Enable}_E$
abbreviates only its local chain: authorization precedes sending its ready value,
and receipt of the peer ready value precedes its final owner CAS and enable event.
No cross-host atomic commit, simultaneous enablement, or global-clock order is
claimed. The effect-producing endpoint must itself be enabled, and its trusted
purpose adapter must commit the unique effect before issuing $r$; receipt syntax
alone implies neither fact.
Equation~\eqref{eq:central-goal} does not claim that a remote endpoint's policy
report or provider execution is hardware-attested.

\section{The Policy-to-Purpose-to-Byte Contract}

\subsection{Policy snapshot and local decision}
A signed policy and monotonic compare-and-swap first yield an immutable policy
snapshot. The trusted service then authenticates a session request and creates a
local decision bound to the caller, local identity and role, expected peer,
exact purpose, exact local owner, protocol version, one wire profile, one local
provider, endpoint-private typed effect scope, policy identity, and decision epoch. No later layer may consult UI
state, environment state, or a second selection policy. The peer-visible decision digest excludes the private
local provider and owner; each endpoint separately proves that its local
decision refines to a compliant implementation of the agreed wire profile. The
handle is a caller-bound, non-serializable reference to service state, not a
forgeable or reflective IPC data object.

\subsection{Closed purpose and canonical context}
The purpose is a closed tagged union. A remote-desktop value binds view-only or
interactive authorization and a canonical capability bitmap. A file-transfer
value binds role-relative direction, transfer identifier, declared length, and
SHA-256 digest. Policy authorizes the exact peer and purpose before creating the
local decision. Generic metadata maps and caller-defined strings are rejected.

\contextv{} is a prefix-free canonical encoding of roles, peer identities,
recipient KEM-key digest, the peer-visible decision projection, full purpose,
client nonce, and a fresh peer-agreed session identifier. BoundSessionV1 has no
suite re-selection and no simultaneous-open subprotocol: an outer arbiter chooses
one initiator and transport owner first. Each endpoint's connection generation
remains private local state. Unknown critical fields, duplicates, non-canonical
order, overflow, and trailing bytes fail before KEM execution.

The pre-KEM context cannot contain a responder nonce or ciphertext that does not
yet exist. Message signatures and the root key schedule therefore bind the later
complete transcript. This two-stage construction avoids a circular definition
while keeping pre-KEM and final-transcript agreement as separate proof
obligations.

\subsection{Grant authorization, bilateral readiness, and receipt}
Message A authenticates the exact context projection and initiator KEM material.
Message B confirms the fixed wire profile and authenticates both identities and nonces, responder
material, and Message A. Direction-specific Finished messages confirm the
complete transcript. A second key schedule separates remote video, input,
file, and control channels by purpose and direction.

After Finished, each local transaction compare-and-swaps its exact private owner
and decision epoch, then atomically stores a pending grant and a
matching authorization record. The record covers the complete local security
tuple and means ``this service authorized this grant,'' not application success.
Each endpoint then sends an authenticated \textsf{GrantReadyV1}. A canonical
shared grant identifier covers the ordered peers, policy, wire profile, peer
session, purpose, context, and transcript; a role-ordered digest covers both
ready messages. Local grant identifiers, decision epochs, owners, and commit
identifiers remain private. Before enabling and publishing \grant{}, the service
again validates its exact owner and stores the bilateral-ready digest. Crash
recovery revalidates and resumes the same pending transition or revokes it.

Business success requires a direction-specific, purpose-keyed receipt binding
the shared grant, bilateral-ready digest, issuer authorization commitment,
purpose, context, transcript, operation, sequence, result, and effect digest.
Only the corresponding sealed purpose committer may report success: ordinary application code
cannot choose the result or digest. Remote presentation means a completed trusted
renderer fence, control means the OS input API accepted the exact action, and
file success requires verified bytes plus the defined durable-store commit.

\subsection{Failure and downgrade}
A policy-bound PQ/T attempt cannot yield a business-capable classic session.
Timeout, network manipulation, overload, cancellation, parse failure, identity
failure, ABI failure, and provider failure terminate explicitly. A separately
identified compatibility protocol may carry an authenticated endpoint
declaration, but that declaration cannot prove the truth of a local hardware
failure. A control-only bootstrap purpose cannot derive business keys before the
required PQ/T rekey.

\section{Layered Security Analysis}

The proof plan deliberately separates claims that rely on different models.
The combiner layer should establish canonical-encoding injectivity and binding
under stated computational assumptions. A new Tamarin theory should model the
actual \system{} messages, signed decision identity, purpose, two-stage context,
Finished exchange, compromises, and application acceptance under an active
network adversary. A second state model should cover policy CAS, crash/restart,
owner replacement, bilateral readiness, policy update, effect journaling, and
publication.

The complete responsibility matrix is maintained with the formal composition
boundary and will move to the supplement once theorem and artifact names are
stable. A check at one assurance layer never discharges another layer.

The minimum theorem set includes policy-snapshot origin and monotonicity, exact
peer-and-purpose authorization by each local decision, local provider refinement
to the agreed wire profile, no policy-bound business fallback, canonical pre-KEM
context agreement, Finished injective agreement, hybrid secrecy with explicit
compromise conditions, no cross-purpose key reuse, exact local-owner transaction,
bilateral-ready agreement, policy-update semantics, trusted effect authority,
authorization-before-publication, and authenticated receipt-to-effect
correspondence. A static KEM
suite must also include an attack witness showing the absence of forward secrecy
after later static KEM-key compromise unless a fresh ephemeral contribution is
specified and proved.

Component binding plus a symbolic protocol proof is not a complete computational
AKE proof. Likewise, vectors, hashes, and source-bound execution are conformance
evidence rather than a mechanized specification-to-binary refinement. These
limits remain in the abstract, threat model, and conclusion after results are
added.

\section{Implementation Plan and Evidence Chain}

The intended implementation has a portable Rust protocol core and narrow ports
for policy, KEM, identity, transactional storage, transport, diagnostic audit,
authenticated IPC, time, and three typed effect committers. Platform bindings
own memory-safe marshalling and I/O. Application modules consume only typed
grants. Existing component implementations may be reused behind ports only after
their revisions, assumptions, and error semantics are pinned; obsolete parallel
protocol paths are removed rather than hidden behind compatibility wrappers.

Every asynchronous boundary revalidates the exact private local owner tuple.
Owner/decision compare-and-swap occurs both in the pending-grant authorization
transaction and immediately before bilateral enable, with no external await in
either transaction. Caller-bound handles,
queues, parser inputs, pre-authentication state, file sizes, operation sequences,
and audit retention have explicit limits, expiry, revocation, and release rules.
Expensive signature and KEM work is admitted only after the cheapest trustworthy
identity and resource checks.

One evidence manifest selects the final normative specification, formal models,
proof reports, source, generated bindings, dependency identities, binaries,
device runs, raw samples, and table generators. The machine-readable claim
ledger distinguishes design, formal, conformance, runtime, physical-device,
performance, artifact, and publication states.

\section{Evaluation Plan}

The final evaluation is organized by four research questions rather than by
component inventory.

\textbf{RQ1: Does the contract prevent cross-policy, cross-purpose, and
downgrade acceptance?} We will execute active attacks that strip or replace
policy, wire profile, purpose, context, key shares, transcript fields,
Finished, shared-grant/ready values, local owner changes, trusted-effect outcomes,
and receipts. Security-event results
will be reported as $x/N$ with Wilson intervals; unreachability claims come only
from a model.

\textbf{RQ2: Which bindings are load-bearing?} Matched ablations remove policy
digest, purpose, public key, ciphertext, attempt declaration, owner check, audit
ordering, either ready message, bilateral-ready digest, shared-grant binding,
trusted-effect authority, effect-journal ordering, and channel separation one at a time. A deliberately unsafe timeout
fallback is labeled an ablation, not a baseline implementation of TLS or QUIC.

\textbf{RQ3: Does one final core execute consistently across substrates?}
Canonical vectors, independent decoders, cross-language differential tests, ABI
fuzzing, sanitizers, and a two-ISA binary constant-time discriminator run from
the final evidence root. At least one real non-Apple endpoint must complete the
same authenticated protocol and application receipt with an Apple endpoint;
cross-compilation and static contract checks do not count as runtime cells.

\textbf{RQ4: What is the application cost?} Matched final-source baselines use
the same application, transport, endpoints, and completion definition. Remote
desktop is the primary performance/QoE workload and reports handshake-to-first-
frame, input-to-authenticated-apply receipt, frame quality, CPU, memory, and
energy. File transfer is retained as an independent cross-purpose correctness,
attack, and durable-receipt workload; its extended performance appears only in
the supplement. Primary non-inferiority margins and experimental units are
declared before capture.

Independent process batches are the experimental unit; frames within one trace
are not independent sessions. Reports bind hardware, OS, compiler, backend,
power/thermal state, transport, exclusions, confidence intervals, and effect
sizes.

\section{Discussion and Limitations}

The proposed contract narrows the gap between cryptographic choice and
application authority. It neither attests remote policy execution nor survives
trusted-boundary compromise. It does not prove provider-failure facts, prevent
all resource exhaustion, or turn a presentation fence into proof of human
perception. Without an external witness, local policy and audit state may roll
back with the device. Stable policy identifiers may create linkability; metadata
privacy requires separate work.

The strongest planned formal result remains layered: computational combiner
binding, symbolic protocol guarantees, and state-machine invariants under
explicit interfaces. A full computational AKE reduction and mechanized
implementation refinement would be stronger and are not implied by the proposed
evidence chain.

\section{Conclusion}

Policy-to-purpose-to-byte continuity makes application acceptance the endpoint
of post-quantum migration assurance. \system{} proposes an independently scoped
contract designed to bind each endpoint's local authenticated policy decision
and exact owner to one wire profile, purpose, transcript, shared grant,
bilateral-ready set, and trusted purpose-specific effect receipt. Its
authorization record is local diagnostic
evidence, not a claim of durable non-repudiation or application success. The next
milestone is not prose expansion: it is a
single final-source implementation and composed evidence root that can replace
the pending entries in the claim ledger without borrowing conclusions from its
upstream components.

\bibliographystyle{IEEEtran}
\bibliography{references}

\end{document}
